%\section{Deception Attacks and Decision-Making}
\section{Decision-Making and Deception Attacks}
\label{sec:decision-making}

\addition{Beyond effects on information channels, we seek to model how \MR deception attacks impact human cognition.} 
% \removal{In this section, we} 
To address \textbf{RQ3}, which concerns the interactions between decision-making and \MR deception attacks, \addition{we} conduct a thorough review of the cognition literature and develop a comprehensive decision-making model that outlines the stages of decision-making susceptible to these attacks.
Figure~\ref{fig:process-model} shows our \MR Decision-Making Model. The model provides an overview of how sensory input is cognitively processed by a user to make decisions and where the different types of attacks affect decision making.
%may harm the initial phases of human cognitive processes, subsequently leading to a chain reaction that impacts perception, decision-making, and responses. 
Our decision-making model comprises of seven components: Sensory Inputs, Attention, Perception, Memory, Decision-Making, Decision Execution, and Responses.

\begin{figure*}[ht!]
\centering
  \includegraphics[width=0.95\textwidth]{figures/decision_making_model.pdf}
  \caption{\MR Deception Decision-Making Model. External stimuli (left) are input to cognitive processes (right). Stimuli are first processed by perception. %, where stimuli are selected, organized, and interpreted. 
  Selective attention manages perception on relevant stimuli. Organized stimuli are stored in working memory. Interpreted stimuli are passed to decision-making, where executive attention manages decisions and their execution.} % Channel attacks target sensory inputs, while processing attacks target cognitive processes.}
 % are captured Channel attacks first have a direct impact on the human sensory system, and processing attacks on the other side have an impact on human cognition. The sensory input is processed through different types of attention. Perception comes next and involves three stages: selection, organization, and interpretation of sensory stimuli. Decision-making is the subsequent stage where choices are made based on perceived information, followed by decision execution, where these decisions are acted upon. The final stage is the response, which can be physiological, behavioral, or cognitive.}
  \label{fig:process-model}
 \vspace{-1ex}
\end{figure*}


\begin{itemize}
\itemsep0em
    \item \emph{Sensory Inputs}: Visual, Auditory, Smell, Taste, and Touch are the five different types of sensory inputs that can be affected by deception attacks.
    \item \emph{Attention}: Initial stage where sensory information is gathered. Provides a gateway to perception.
    \item \emph{Perception}: Sensory information gathered is processed and understood.
    \item \emph{Memory}: Processed information is stored in working or long-term memory for future use and retrieval.
    \item \emph{Decision-Making}: Determining a particular course of action predicated on perception.
    \item \emph{Decision Execution}: Decisions are executed.
    \item \emph{Responses}: Physiological, behavioral, or cognitive responses of executed decisions. %Handled by motor processor.
\end{itemize}

\subsection{Perception}
Perception refers to the cognitive process through which one comprehends sensory stimuli~\cite{qiong2017brief}.
Wang et al.~\cite{wang2007cognitive} define perception as ``a set of internal sensational cognitive processes of the brain at the subconscious cognitive function layer that detects, relates, interprets, and searches internal cognitive information in the mind.'' 
Perception is either active or passive.
Active perception involves the intentional direction of attention towards environmental stimuli to extract information \cite{gibson2014ecological}. 
In contrast, passive perception occurs without deliberate effort; sensory information is received as presented \cite{Rock1983-ROCTLO}.

% Gibson \cite{gibson2002theory} uses the term perception for any experience of the environment surrounding the body and the term proprioception or internal perception for any experience of the body itself. One example of perception is the ability of vision to identify an object's motion in relation to a motionless environment. Similarly, one example of proprioception is one's movement in relation to a fixed environment, whether active or passive can be detected by vision.

Perception involves three stages~\cite{qiong2017brief}:
\begin{itemize}
\itemsep0em 
    \item \emph{Selection:} Filter and select environmental stimuli from meaningful experiences.
    \item \emph{Organization:} Structure and categorize the selected information, creating coherent and stable perceptions through grouping by proximity and similarity.
    \item \emph{Interpretation:} Assign meaning to organized stimuli, with individuals' cultural or experiential backgrounds leading to different understandings of the same stimuli.
\end{itemize}

In each stage of perception, \MR deception attacks can target specific vulnerabilities. 
%Dionisio et al. \cite{Dionisio2001Differentiation} highlight the cognitive demands associated with generating deceptive responses, emphasizing the increased cognitive load during these processes. 
%Their findings suggest that deception in \MR not only challenges the perception stages but also significantly taxes the cognitive resources required for maintaining accurate perception.
During selection, attacks can cause sensory overload or misdirect focus on irrelevant stimuli. In the organization step, attacks could involve boundary or dimension manipulation, affecting how stimulus are structured and grouped due to changes in proximity or scale. Propaganda or spoofing attacks can target interpretation, affecting the meaning assigned to stimuli that may seem wrong but is coming from a trusted source (e.g., the system or a collaborator). These potential attacks highlight the importance of the accuracy and reliability of perception in \MR systems.

In addition to the conscious components of perception, subliminal inputs play an important role in how individuals interact with and understand \MR environments.
Cetnarski et al.\ \cite{Cetnarski2014Subliminal} show that subliminal stimuli—information presented below the threshold of conscious awareness—can significantly influence decision-making processes in \MR. This underscores the need to understand these subtle interactions that occur at the subconscious level of perception.
% Gestalt Principles of Perception 
% Perception is the main factor in defining individuals' interpretation of their environment. The accuracy and reliability of perception become critical in \MR systems, where the boundary between the real and virtual world often becomes indistinct.
%Deception attacks can exploit this by altering sensory inputs or manipulating the processing of these inputs, leading to users' misunderstanding. In \MR systems, perception is tightly connected to users' trust. When sensory inputs are reliable and the virtual environment aligns with reality, trust in the system improves. When deception attacks alter perception, mistrust in the system increases, and users may question their interaction with the system. Risk perception, also, is heavily related to trust. As users obtain an increased awareness of the reliability of their perception in an \MR environment, they might perceive increased risks in their actions. Factors like the gravity of events, media coverage, and individual attributes such as age, gender, and each person's past experiences influence how risks are perceived and subsequently, how trust in the system is established or damaged. 
%As mentioned in the previous section, there are two types of attacks on \MR systems, Channel Attacks and Processing Attacks. In our general decision-making model, these two types of attacks are interrelated with specific components of decision-making defined based on where they hold a higher impact in the cognitive decision-making process.


% \subsection{Sensory Input}

% Sensory inputs are frequently the primary focus of Channel Attacks, as their primary objective lies in heightening ambiguity and provoking misleading perceptions; This is achieved through the manipulation and alteration of users' conduct throughout their interactions with the system. For instance, Denial Attacks, have the effect of limiting the range of perceptible information, thereby forcing users to make decisions based on incomplete data. This limited input might also shift users' attention to look for the reasons behind the absence or unavailability of such information. Such situations can be exploited by attackers to divert attention from more critical areas of the system.
% Overt Degradation attacks are apparent and noticeable, making users immediately aware of the compromise in sensory input quality. Sensory overload is a vivid example of these attacks. By overwhelming users with excessive stimuli, these attacks have the potential to affect rational thinking and limit decision-making ability, which can lead to immediate distrust and a reevaluation of the system's reliability.
% Covert degradation attacks, in contrast, subtly degrade the quality of sensory inputs over time, making them harder to detect. The danger with covert attacks lies in their subtle and deceptive characteristics, which may result in users failing to immediately recognize them, leading them to make decisions based on compromised data without realizing it.
%the Human Joystick Attack has the potential to significantly impact users' tactile perception, while the Overlay attack could instead disrupt users' visual experience. This impact has a significant effect on users' decision-making abilities and alters their decision responses by interfering with and influencing their sensory inputs.

%\vspace{-1ex}
\subsection{Attention} 

James~\cite{James1890} described attention as the cognitive process by which the mind selectively concentrates on a singular element from a variety of possible stimuli or thoughts, emphasizing its essential function in creating our conscious perception. 
The seminal work of Posner \cite{posner1980orienting} 
%further delineated attention into more specific types, reflecting its multifaceted nature. Posner
introduced a framework for understanding the neural bases of attention and its various components and extending James's initial descriptions into a more nuanced understanding of the brain's attentional mechanisms. Building upon these early foundations, attention classification has expanded to include four types: %Selective, Divided, Sustained, and Executive.

% From these early foundations, the classification of attention has evolved to include several key types: Selective, Divided, Sustained, and Executive.

\begin{itemize}
    \itemsep0em
    \item \emph{Selective:} Focusing on relevant information while suppressing irrelevant information \cite{stevens2012role, murphy2016twenty}.
    
    \item \emph{Divided:} Capacity to allocate cognitive resources to multiple stimuli simultaneously, enabling individuals to engage in concurrent activities~\cite{spelke1976skills}. 
    %In other words, it refers to situations in which two or more channels are attended simultaneously. 
    Attended stimuli are from the same sensory modality or different ones~\cite{Herbranson2017}.

    \item \emph{Sustained:} Readiness to perceive and respond to stimuli over prolonged periods, often without conscious awareness of this vigilance~\cite{mackworth1948breakdown}.
    % Sustained attention, as outlined in the seminal work by Mackworth \cite{mackworth1948breakdown},refers to the psychological readiness to perceive and respond to stimuli over prolonged periods without necessarily being conscious of this vigilance.
    % Degradation attacks may limit sustained attention by making it more challenging for the user to maintain their focus over time, particularly when the quality of sensory inputs fluctuates or declines, resulting in increased cognitive load.
    \item \emph{Executive:} Regulates cognitive and emotional responses through management of other cognitive processes~\cite{posner1990attention}. Aids orchestration of thought and emotion in alignment with goals and the dynamic demands of the environment.
    % Subversion attacks could challenge executive attention by forcing users to constantly adapt to unexpected or counterintuitive system responses, requiring continuous updating of working memory.
\end{itemize}

Channel attacks primarily target Selective and Sustained attention. 
They manipulate the sensory channels through which users receive information, affecting their ability to focus on relevant stimuli or maintain attention over time.
Selective attention is exploited by degrading the sensory inputs, making it harder for users to distinguish between relevant and irrelevant stimuli. As mentioned in Section \ref{sec:decattacks}, this happens in overt and covert degradation attacks. These attacks may also limit sustained attention by making it more challenging for the user to maintain their focus over time, particularly when the quality of sensory inputs fluctuates or declines, resulting in increased cognitive load. 
Denial attacks block access to certain stimuli or information channels, disrupting selective attention. 
%This may potentially cause cognitive overload to select the relevant stimuli.
%A classic framework for understanding this phenomenon is provided by Broadbent's Filter Model of selective attention, which describes how individuals filter information based on physical characteristics before processing its meaning \cite{broadbent2013perception}.

Processing attacks primarily affect Divided and Executive attention by overloading the cognitive processing capabilities or by requiring constant adjustments to unexpected system behaviors.
Corruption attacks %involve the deliberate manipulation or alteration of sensory data to mislead users.
can directly impact the users' selective and executive attention by altering the information presented within \MR environment and also exploiting perceptual biases.
Subversion attacks could challenge executive attention by forcing users to constantly adapt to unexpected system responses, requiring continuous updating of working memory. They also can target divided attention by interrupting the flow of tasks or actions within an \MR environment, which compels users to divide their attention between correcting system errors and accomplishing their original goals.

\finding{Perception and attention are the primary targets for \MR deception attacks. Channel attacks target selection mechanisms by degrading or denying stimuli. Processing attacks target interpretation and execution by corrupting beliefs or subverting interpretations.}
\vspace{-1ex}

%Users' Attention can be manipulated or misconstrued by the employment of processing attacks; These attacks aim to effectively manipulate user perception, leading to false renditions and inferences. Incorrect or skewed perception then leads to misguided decision-making, which then impacts decision execution and the various types of responses.
%Corruption and Subversion attacks are specifically aimed at undermining the fundamental mechanisms that control attention within \MR systems.
%Corruption attacks have the potential to introduce anomalies or unwanted modifications in sensory data, drawing users' attention to these altered elements. This change in attention can potentially result in users failing to perceive or acknowledge other essential data, leading to a distorted understanding of their environment. Similarly, subversion attacks, by making the system behave in unintended ways, can create distractions or unexpected events that capture users' attention.
%The manipulation of attention through these attacks has particularly negative impacts on the \MR system. When attention is diverted or overwhelmed, the quality and accuracy of perception are reduced. Users might interpret the elements of the environment incorrectly. This compromised perception, affects decision-making, as decisions are based on this inaccurate understanding. As a result of this, decision execution might not be right or even dangerous, especially whenever immediate responses are needed.


\definecolor{lowcolor}{HTML}{bcabcd} % Light red
\definecolor{midcolor}{HTML}{8a6ca8} % Light yellow
\definecolor{highcolor}{HTML}{582D83}
\definecolor{none}{HTML}{ffffff}


% \begin{table}[htbp]
%     \centering
%     \begin{tabular}{|l|c|}
%         \hline
%         \rowcolor{gray!30}
%         \textbf{Category} & \textbf{Value} \\
%         \hline
%         Low & \cellcolor{green!30} Low \\
%         Medium & \cellcolor{yellow!30} Medium \\
%         High & \cellcolor{red!30} High \\
%         \hline
%     \end{tabular}
%     \caption{Visualization of Low, Medium, and High Values}
%     \label{tab:visualization}
% \end{table}


\newcommand{\Low}{%
\begin{tikzpicture}
\fill[fill=lowcolor] (0,0) rectangle (0.33,0.2);
\fill[fill=none] (0.33,0) rectangle (0.66,0.2);
\fill[fill=none] (0.66,0) rectangle (1,0.2);
\end{tikzpicture}%
}

\newcommand{\Medium}{%
\begin{tikzpicture}
\fill[fill=midcolor] (0,0) rectangle (0.33,0.2);
\fill[fill=midcolor] (0.33,0) rectangle (0.66,0.2);
\fill[fill=none] (0.66,0) rectangle (1,0.2);
\end{tikzpicture}%
}

\newcommand{\High}{%
\begin{tikzpicture}
\fill[fill=highcolor] (0,0) rectangle (1,0.2); % Use custom color
%\fill[fill=highcolor] (0.33,0) rectangle (0.66,0.2);
%\fill[fill=highcolor] (0.66,0) rectangle (1,0.2);
\end{tikzpicture}%
}

\newcommand{\LowMedium}{%
\begin{tikzpicture}
\fill[fill=lowcolor] (0,0) rectangle (0.33,0.2);
\fill[fill=midcolor] (0.33,0) rectangle (0.66,0.2);
\fill[fill=none] (0.66,0) rectangle (1,0.2);
\end{tikzpicture}%
}

\newcommand{\MediumHigh}{%
\begin{tikzpicture}
\fill[fill=none] (0,0) rectangle (0.33,0.2);
\fill[fill=midcolor] (0.33,0) rectangle (0.66,0.2);
\fill[fill=highcolor] (0.66,0) rectangle (1,0.2);
\end{tikzpicture}%
}

\newcommand{\LowHigh}{%
\begin{tikzpicture}
%\fill[left color=lowcolor, right color=highcolor] (0,0) rectangle (1,0.2);
\fill[fill=lowcolor] (0,0) rectangle (0.33,0.2);
\fill[fill=midcolor] (0.33,0) rectangle (0.66,0.2);
\fill[fill=highcolor] (0.66,0) rectangle (1,0.2);
\end{tikzpicture}%
}



\begin{table*}[ht!]
\caption{\MR attacks from our ontology are assessed according to the Information-Theoretic Model and Decision-Making Model.}
%mapped to the type of attack and the Information-Theoretic Model. Then, these attacks are assessed according to the degree to which they affect perception, attention, and memory. }





% Taking into account our \MR Deception Information-Theory and Decision-Making models, each attack is evaluated based on how it affects information-theoretic characteristics to alter communication channels and cognitive processes to determine the extent of its interference with various stages of human cognition.

\small
\centering
\renewcommand{\arraystretch}{1.15}
\setlength{\tabcolsep}{1pt}

\begin{tabular}{ccr|cccc|cccccccc|c|}
\cline{4-16}
  & &  & 
  \multicolumn{4}{c|}{Information-} & 
  \multicolumn{9}{c|}{} \\   % First row for spacing

  
  & &  & 
  \multicolumn{4}{c|}{\multirow{-1}{*}{Theoretic Model}} & 
  \multicolumn{9}{c|}{\multirow{-2}{*}{Decision-Making Model}} \\ \cline{4-16} % Second row, "Information Theoretic Model"

  
  & &  & 
  \multicolumn{3}{c|}{$C$} & & 
  \multicolumn{4}{c|}{Perception} & 
  \multicolumn{4}{c|}{Attention} & 
  \multirow{2}{*}{Mem.} \\ \cline{4-6} \cline{8-15} % Third row, details

  
  &
   &
  \multicolumn{1}{r|}{\multirow{-2}{*}{ \backslashbox{Attacks}{Models}}} &
  \multicolumn{1}{c|}{\hspace*{1mm}$W$\hspace*{1mm}} &
  \multicolumn{1}{c|}{\hspace*{1mm}$S$\hspace*{1mm}} &
  \multicolumn{1}{c|}{\hspace*{1mm}$N$\hspace*{1mm}} &
  \multirow{-2}{*}{$M$} &
  \multicolumn{1}{c|}{\hspace*{0.5mm}A/P\hspace*{0.5mm}} &
  \multicolumn{1}{c|}{$Sel$} &
  \multicolumn{1}{c|}{$Org$} &
  \multicolumn{1}{c|}{$Int$} &
  \multicolumn{1}{c|}{$Foc$} &
  \multicolumn{1}{c|}{$Div$} &
  \multicolumn{1}{c|}{$Sus$} &
  \multicolumn{1}{c|}{$Exe$} &
  \multicolumn{1}{c|}{} \\ \hlinewd{1.5pt}

  \multicolumn{1}{|c|}{} & \multicolumn{1}{c|}{}
   &
  Sensory Overload~ &
  \multicolumn{1}{c|}{} &
  \multicolumn{1}{c|}{} &
  \multicolumn{1}{c|}{\ding{51}} &
   &
  \multicolumn{1}{c|}{A} &
  \multicolumn{1}{c|}{\High} &
  \multicolumn{1}{c|}{\High} &
  \multicolumn{1}{c|}{\High} &
  \multicolumn{1}{c|}{\High} &
  \multicolumn{1}{c|}{\High} &
  \multicolumn{1}{c|}{\High} &
  \multicolumn{1}{c|}{\High} &
  \multicolumn{1}{c|}{W} \\ 
   \multicolumn{1}{|c|}{} & \multicolumn{1}{c|}{}
   &
  Momentary Misdirection~ & 
  \multicolumn{1}{c|}{} &
  \multicolumn{1}{c|}{} &
  \multicolumn{1}{c|}{\ding{51}} &
   &
  \multicolumn{1}{c|}{A} &
  \multicolumn{1}{c|}{\LowHigh} &
  \multicolumn{1}{c|}{\LowHigh} &
  \multicolumn{1}{c|}{\LowHigh} &
  \multicolumn{1}{c|}{\High} &
  \multicolumn{1}{c|}{\LowHigh} &
  \multicolumn{1}{c|}{\LowHigh} &
  \multicolumn{1}{c|}{\LowHigh} &
  \multicolumn{1}{c|}{W} \\ 
  \multicolumn{1}{|c|}{} &
  \multicolumn{1}{c|}{\multirow{-3}{*}{Overt} }&
  Signal Replacement~ &
  \multicolumn{1}{c|}{} &
  \multicolumn{1}{c|}{\ding{51}} &
  \multicolumn{1}{c|}{} &
   &
  \multicolumn{1}{c|}{A} &
  \multicolumn{1}{c|}{\High} &
  \multicolumn{1}{c|}{\High} &
  \multicolumn{1}{c|}{\High} &
  \multicolumn{1}{c|}{\High} &
  \multicolumn{1}{c|}{\High} &
  \multicolumn{1}{c|}{\High} &
  \multicolumn{1}{c|}{\High} &
  \multicolumn{1}{c|}{W} \\  
 \multicolumn{1}{|c|}{} & 
  \multicolumn{1}{c|}{\multirow{-3}{*}{Degradation}} & 
  Quality Erosion~ &
  \multicolumn{1}{c|}{\ding{51}} &
  \multicolumn{1}{c|}{} &
  \multicolumn{1}{c|}{\ding{51}} &
   &
  \multicolumn{1}{c|}{A} &
  \multicolumn{1}{c|}{\LowHigh} &
  \multicolumn{1}{c|}{\LowHigh} &
  \multicolumn{1}{c|}{\LowHigh} &
  \multicolumn{1}{c|}{\LowHigh} &
  \multicolumn{1}{c|}{\LowHigh} &
  \multicolumn{1}{c|}{\LowHigh} &
  \multicolumn{1}{c|}{\LowHigh} &
  \multicolumn{1}{c|}{W} \\ \cline{2-16} 
   \multicolumn{1}{|c|}{} & \multicolumn{1}{c|}{} &
  Subtle Boundary Manipulation~  &
  \multicolumn{1}{c|}{} &
  \multicolumn{1}{c|}{\ding{51}} &
  \multicolumn{1}{c|}{} &
   &
  \multicolumn{1}{c|}{P} &
  \multicolumn{1}{c|}{\Low} &
  \multicolumn{1}{c|}{\Low} &
  \multicolumn{1}{c|}{\LowHigh} &
  \multicolumn{1}{c|}{\Low} &
  \multicolumn{1}{c|}{\Low} &
  \multicolumn{1}{c|}{\Low} &
  \multicolumn{1}{c|}{\LowMedium} &
  \multicolumn{1}{c|}{W} \\ 
 \multicolumn{1}{|c|}{} &
  \multicolumn{1}{c|}{\multirow{-2}{*}{Covert}} &
  Subtle Dimension Manipulation~  &
  \multicolumn{1}{c|}{} &
  \multicolumn{1}{c|}{\ding{51}} &
  \multicolumn{1}{c|}{} &
   &
  \multicolumn{1}{c|}{P} &
  \multicolumn{1}{c|}{\Low} &
  \multicolumn{1}{c|}{\LowMedium} &
  \multicolumn{1}{c|}{\LowHigh} &
  \multicolumn{1}{c|}{\Low} &
  \multicolumn{1}{c|}{\Low} &
  \multicolumn{1}{c|}{\Low} &
  \multicolumn{1}{c|}{\LowMedium} &
  \multicolumn{1}{c|}{W} \\
 \multicolumn{1}{|c|}{} &
  \multicolumn{1}{c|}{\multirow{-2}{*}{Degradation}} &
  Physical Movement Manipulation~ &
  \multicolumn{1}{c|}{} &
  \multicolumn{1}{c|}{\ding{51}} &
  \multicolumn{1}{c|}{} &
   &
  \multicolumn{1}{c|}{P} &
  \multicolumn{1}{c|}{\Low} &
  \multicolumn{1}{c|}{\Low} &
  \multicolumn{1}{c|}{\LowHigh} &
  \multicolumn{1}{c|}{\Low} &
  \multicolumn{1}{c|}{\Low} &
  \multicolumn{1}{c|}{\Low} &
  \multicolumn{1}{c|}{\LowHigh} &
  \multicolumn{1}{c|}{W} \\ \cline{2-16} 
 \multicolumn{1}{|c|}{} & \multicolumn{1}{c|}{}
   &
  Shutdown~ &
  \multicolumn{1}{c|}{\ding{51}} &
  \multicolumn{1}{c|}{} &
  \multicolumn{1}{c|}{} &
   &
  \multicolumn{1}{c|}{A} &
  \multicolumn{1}{c|}{\High} &
  \multicolumn{1}{c|}{\High} &
  \multicolumn{1}{c|}{\High} &
  \multicolumn{1}{c|}{\High} &
  \multicolumn{1}{c|}{\High} &
  \multicolumn{1}{c|}{\High} &
  \multicolumn{1}{c|}{\High} &
  \multicolumn{1}{c|}{W} \\ 
  \multicolumn{1}{|c|}{} & \multicolumn{1}{c|}{}
   &
  Overlay~ &
  \multicolumn{1}{c|}{\ding{51}} &
  \multicolumn{1}{c|}{} &
  \multicolumn{1}{c|}{\ding{51}} &
   &
  \multicolumn{1}{c|}{A} &
  \multicolumn{1}{c|}{\LowHigh} &
  \multicolumn{1}{c|}{\LowHigh} &
  \multicolumn{1}{c|}{\LowHigh} &
  \multicolumn{1}{c|}{\LowHigh} &
  \multicolumn{1}{c|}{\LowHigh} &
  \multicolumn{1}{c|}{\LowHigh} &
  \multicolumn{1}{c|}{\LowHigh} &
  \multicolumn{1}{c|}{W} \\  
\multicolumn{1}{|c|}{\multirow{-10}{*}{\rotatebox[origin=c]{90}{Channel Attacks}}} &
  \multicolumn{1}{c|}{\multirow{-3}{*}{Denial}} & 
  Removal~ &
  \multicolumn{1}{c|}{\ding{51}} &
  \multicolumn{1}{c|}{\ding{51}} &
  \multicolumn{1}{c|}{} &
   &
  \multicolumn{1}{c|}{A} &
  \multicolumn{1}{c|}{\LowHigh} &
  \multicolumn{1}{c|}{\LowHigh} &
  \multicolumn{1}{c|}{\LowHigh} &
  \multicolumn{1}{c|}{\LowHigh} &
  \multicolumn{1}{c|}{\LowHigh} &
  \multicolumn{1}{c|}{\LowHigh} &
  \multicolumn{1}{c|}{\LowHigh} &
  \multicolumn{1}{c|}{W} \\ \cline{1-16}
 \multicolumn{1}{|c|}{} &
  \multicolumn{1}{c|}{} &
  Spoofing~ &
  \multicolumn{1}{c|}{} &
  \multicolumn{1}{c|}{} &
  \multicolumn{1}{c|}{} &
  \ding{51} &
  \multicolumn{1}{c|}{P} &
  \multicolumn{1}{c|}{\Low} &
  \multicolumn{1}{c|}{\Low} &
  \multicolumn{1}{c|}{\LowHigh} &
  \multicolumn{1}{c|}{\LowHigh} &
  \multicolumn{1}{c|}{\LowHigh} &
  \multicolumn{1}{c|}{\LowHigh} &
  \multicolumn{1}{c|}{\LowHigh} &
  \multicolumn{1}{c|}{W} \\ 
 \multicolumn{1}{|c|}{} &
  
  \multicolumn{1}{c|}{\multirow{-2}{*}{Corruption}} &
  False-Flag Operations~ &
  \multicolumn{1}{c|}{} &
  \multicolumn{1}{c|}{} &
  \multicolumn{1}{c|}{} &
  \ding{51} &
  \multicolumn{1}{c|}{P} &
\multicolumn{1}{c|}{\Low} &
  \multicolumn{1}{c|}{\Low} &
  \multicolumn{1}{c|}{\LowHigh} &
  \multicolumn{1}{c|}{\Low} &
  \multicolumn{1}{c|}{\Low} &
  \multicolumn{1}{c|}{\Low} &
  \multicolumn{1}{c|}{\Low} &
  \multicolumn{1}{c|}{W/L} \\ \cline{2-16} 
  \multicolumn{1}{|c|}{} &
  \multicolumn{1}{c|}{} &
  Bias Attacks~ &
  \multicolumn{1}{c|}{} &
  \multicolumn{1}{c|}{} &
  \multicolumn{1}{c|}{} &
   &
  \multicolumn{1}{c|}{P} &
  \multicolumn{1}{c|}{\Low} &
  \multicolumn{1}{c|}{\Low} &
  \multicolumn{1}{c|}{\LowHigh} &
  \multicolumn{1}{c|}{\Low} &
  \multicolumn{1}{c|}{\Low} &
  \multicolumn{1}{c|}{\Low} &
  \multicolumn{1}{c|}{\Low} &
  \multicolumn{1}{c|}{W/L} \\  
  \multicolumn{1}{|c|}{} &
  \multicolumn{1}{c|}{} &
  Lure~ &
  \multicolumn{1}{c|}{} &
  \multicolumn{1}{c|}{} &
  \multicolumn{1}{c|}{} &
   &
  \multicolumn{1}{c|}{P} &
  \multicolumn{1}{c|}{\LowHigh} &
  \multicolumn{1}{c|}{\LowHigh} &
  \multicolumn{1}{c|}{\LowHigh} &
  \multicolumn{1}{c|}{\LowHigh} &
  \multicolumn{1}{c|}{\LowHigh} &
  \multicolumn{1}{c|}{\LowHigh} &
  \multicolumn{1}{c|}{\LowHigh} &
  \multicolumn{1}{c|}{W/L} \\  
  \multicolumn{1}{|c|}{} &
  \multicolumn{1}{c|}{} &
  Disinformation~ &
  \multicolumn{1}{c|}{} &
  \multicolumn{1}{c|}{} &
  \multicolumn{1}{c|}{} &
   &
  \multicolumn{1}{c|}{P} &
  \multicolumn{1}{c|}{\Low} &
  \multicolumn{1}{c|}{\Low} &
  \multicolumn{1}{c|}{\LowHigh} &
  \multicolumn{1}{c|}{\Low} &
  \multicolumn{1}{c|}{\Low} &
  \multicolumn{1}{c|}{\Low} &
  \multicolumn{1}{c|}{\Low} &
  \multicolumn{1}{c|}{W/L} \\ 
  \multicolumn{1}{|c|}{} &
  \multicolumn{1}{c|}{} &
  Propaganda~ &
  \multicolumn{1}{c|}{} &
  \multicolumn{1}{c|}{} &
  \multicolumn{1}{c|}{} &
   &
  \multicolumn{1}{c|}{P} &
  \multicolumn{1}{c|}{\Low} &
  \multicolumn{1}{c|}{\Low} &
  \multicolumn{1}{c|}{\LowHigh} &
  \multicolumn{1}{c|}{\Low} &
  \multicolumn{1}{c|}{\Low} &
  \multicolumn{1}{c|}{\Low} &
  \multicolumn{1}{c|}{\Low} &
  \multicolumn{1}{c|}{W/L} \\ 
\multicolumn{1}{|c|}{\multirow{-7}{*}{\rotatebox[origin=c]{90}{Processing Attacks}}} &
  \multicolumn{1}{c|}{\multirow{-5}{*}{Subversion}} &
  Gaslighting~ &
  \multicolumn{1}{c|}{} &
  \multicolumn{1}{c|}{} &
  \multicolumn{1}{c|}{} &
   &
  \multicolumn{1}{c|}{P} &
  \multicolumn{1}{c|}{\Low} &
  \multicolumn{1}{c|}{\Low} &
  \multicolumn{1}{c|}{\LowHigh} &
  \multicolumn{1}{c|}{\Low} &
  \multicolumn{1}{c|}{\Low} &
  \multicolumn{1}{c|}{\Low} &
  \multicolumn{1}{c|}{\LowHigh} &
  \multicolumn{1}{c|}{W/L} \\ \hline
\end{tabular}
\\ \vspace{0.5em}
Low = 
\begin{tikzpicture}
\fill[fill=lowcolor] (0,0) rectangle (0.33,0.2); 
\end{tikzpicture}, 
Low-Medium = \begin{tikzpicture}
\fill[fill=lowcolor] (0,0) rectangle (0.33,0.2);
\fill[fill=midcolor] (0.33,0) rectangle (0.66,0.2);
\end{tikzpicture}, Low-High = \LowHigh, High = \High \protect\\
  \textbf{Information-Theoretic Model}: $C =$ Channel Capacity, $W =$ Bandwidth, $S =$ Signal, $N =$ Noise, $M =$ Mimicry \protect\\ \textbf{Perception}: $A/P =$ Active/Passive, $Sel =$ Selection, $Org =$ Organization, $Int =$ Interpretation \protect\\ \textbf{Attention}: $Foc =$ Selective, $Div =$ Divided, $Sus =$ Sustained, $Exe =$ Executive; \textbf{Memory}: $W =$ Working, $L =$ Long-Term  
\label{tab:framework}
\vspace{-1ex}
\end{table*}


\subsection{Memory}
Working memory and long-term memory are central components of our decision-making model. 
Baddeley~\cite{baddeley1974working,alan1992working,baddeley2007working} derived a multicomponent model of working memory consisting of the visuospatial sketchpad, phonological loop, central executive, and episodic buffer.
The visuospatial sketchpad stores visual and spatial information while the phonological loop stores auditory and verbal information.
The central executive directs attention towards \minoraddition{stored} information \minoraddition{in either one}.
% \minorremoval{stored in either the visuospatial sketchpad or the phonological loop }
The episodic buffer provides temporary storage of information needed by the central executive with connections to the other three components and long-term memory.
Long-term memory represents a permanent store that receives selected inputs from both the sensory register and working memory~\cite{ATKINSON196889}.


%Deception in \MR systems challenges users' perceptions of virtual and real-world content, 
%increasing cognitive load \cite{Dionisio2001Differentiation} and 
%affecting memory recall and decision-making \cite{Sporer2016Deception}.%, which is critical for recalling specific experiences. 
% Working memory and attention enable goal-driven processing by increasing access to relevant information \cite{awh2006interactions}. 
% Working memory influences the deployment of attention to new information and the manipulation and updating of existing representations.
% Attention affects the encoding and maintenance of information in working memory during multiple stages of processing, from early sensory to post-perceptual stages. 


\MR deception attacks affect memory and correspondingly attention.
Downing \cite{downing2000interactions} showed that the content of visuospatial sketchpad can guide selective attention toward matching visual stimuli. 
Through spoofing attacks, adversaries can produce deceptive stimuli that match expected stimuli, leveraging working memory to direct the user's selective attention.
Santangelo and Macaluso~\cite{santangelo2013contribution} identified the critical role of working memory in managing divided attention
%, underscoring the substantial cognitive demands placed on individuals 
when monitoring multiple objects simultaneously. 
Working memory load directly affects the efficiency of the central executive, with increased load impairing attention to multiple stimuli. 
Thus, sensory overload attacks can overwhelm working memory by visualizing too many objects for working memory to maintain.
Unsworth \& Robinson \cite{unsworth2020working} suggested 
%that individual differences in \WMC significantly impact sustained attention. They indicate 
that individuals with lower \WMC may struggle more with maintaining consistent attention, leading to performance degradation in tasks requiring prolonged focus. 
Therefore, the impact of \MR deception attacks that target \WMC, such as sensory overload, will vary from person to person.
%Engle \cite{engle2002working} suggested that \WMC is about controlling attention, rather than merely the quantity of information that can be stored. 
%This capacity enables individuals to maintain information in an easily accessible state while resisting interference, which is crucial for performing various cognitive tasks, including decision-making.

%Deception, characterized by the intentional manipulation of perception, has a profound impact on cognitive implications. 
%Based on a study by Sporer~\cite{Sporer2004The}, the challenge lies in differentiating between true and deceptive memory cues, which is intricately linked to reality monitoring. Trivers~\cite{Trivers2011Deceit} further elaborates on the intricate relationship between deceit and self-deception, emphasizing their cognitive cost and potential benefits in various interactions.

%Deception increases cognitive load \cite{Dionisio2001Differentiation}, affecting memory recall and decision-making \cite{Sporer2016Deception}. 
%Dionisio et al.~\cite{Dionisio2001Differentiation} found that generating deceptive responses requires greater cognitive processing, as indicated by increased pupillary responses. 
%Sporer~\cite{Sporer2016Deception} suggests the need for a broader understanding of cognitive processes in deception and a more comprehensive framework. 
%Cranford et al.'s~\cite{Cranford2021Towards} concept of cyber deception further exemplifies this need, especially in \MR environments where traditional physical cues are absent or altered.
% Deception activates specific brain regions, differentiating between true and false memories.
% Abe et al. \cite{abe2018neural} found that deception and false memory recruit different brain regions, such as the left temporoparietal cortex for true memory and the right anterior hippocampus for false memory. They identified the neural correlates of true memory, false memory, and deception and found that the left temporoparietal regions were more active for true memory, whereas the right anterior hippocampus was more active for false memory.
% Lee et al. \cite{Lee2009Are} demonstrate that intentionally faked responses and unintentional errors activate different areas of the brain, including the left ventrolateral prefrontal region (BA 47), the right posterior cingulate region (BA 23), and the left precuneus. 

% Cetnarski et al. \cite{Cetnarski2014Subliminal} show that subliminal stimuli in \MR can significantly influence decision-making, underscoring the need to understand these subtle interactions.





% \subsection{Sensemaking}

% Both Channel and Processing attacks impact sensemaking, which relies on processing information presented through \MR headsets.
% Channel attacks prevent the effective collecting and organizing of information in sensemaking tasks.
% Denial attacks deny access to certain information, preventing \MR users from adequately making sense of a situation due to lack of information.
% Overt Degradation attacks disrupt an \MR user's ability to engage in sensemaking by overloading senses with too much information to process when sensemaking.
% Covert Degradation attacks seek to mislead sensemaking by subtly reducing the quality of information hindering how \MR users make sense of a situation.
% For example, an attacker may lower the brightness of holographic callouts that label engine parts with abnormal readings causing a mechanic to miss information that may be valuable to making sense of the problem. 
% Processing attacks involve altering or fabricating the information presented through \MR and can profoundly affect the sensemaking process by presenting users with distorted or completely false information. 
% In Corruption attacks, the attacker might alter the data shown in an \MR environment, which can lead users to incorrect interpretations and misguided actions, since they base their decisions on falsified information.
% Subversion attacks are about covertly manipulating the system to misrepresent data, which might involve altering the way virtual objects behave or interact in the \MR environment, leading to confusion and misjudgment.


%Corruption attacks ...
%Subversion attacks ...
 

%\todo[inline]{explain how different types of attacks affect cognitive components}

%\todo[inline]{how is perception affected? discuss trust and risk perception factors}